% Transcription — fonds Alexandre Grothendieck, shelfmark 107, batch 2 (pages 21-25).
% Established from the facsimile archives/batches/107/batch-02.pdf, cut from
% a copy of 107.pdf fetched through the project's own relay
% (https://grothendieck-relay.onrender.com/source/107.pdf) — not uploaded by
% the user, and not mirrored with `npm run archive`, whose host is blocked —
% per the skill transcribe-grothendieck. The batch file carries no cover
% sheet: PDF page k is folder page 20+k. The folder is twenty-five pages;
% this is the second and last batch, only five pages. Batch 1 (pages 1-20)
% will be transcribed separately and did not exist when this pass was made,
% so the notation below was fixed from these pages alone.
%
% Pass: Opus 5.5 (claude-opus-5-5), 2026-09-26 — first pass, unchecked against the pages by a human.
%
% Pages skipped: 25, the second page of an assembler listing (« ASM 0201
% 15.27 06/04/82 ») with nothing of his mathematics on it. The gap in the
% numbering is the record.
%
% Pages transcribed: 21-24. All four are sheets of continuous-feed computer
% paper. Pages 21 and 23 are printed listings (a JCL job, its allocation
% messages) on which he drew commutative diagrams in pencil, with no prose;
% the listing is not transcribed, the diagrams are. Page 23's listing is
% stamped 82155 (day 155 of 1982, i.e. early June 1982), which fits the
% archivists' « [à partir de 1982] »; this is said once, in a \note, where
% it is observed. Pages 22 and 24 are the blank backs of the same paper,
% written in ink, in English, fast drafting; the formulas carry, a few
% connective words do not.
%
% His pagination: page 22 carries a circled « e » at the top, page 24 a
% circled « d ». The archival order therefore runs e before d; this is kept
% and said in a note at each page. Letters a-c, if they exist, would be in
% batch 1. Page 24 opens « Returning to our initial pb », a reference back
% to pages not in this batch.
%
% Pages summarised: none.
%
% What the batch is about. Closed model categories in the sense of Quillen,
% in the setting of Pursuing Stacks (1983): lifting of homotopies against a
% cofibration i : X -> X' and a fibration p : Y -> Y'. Page 21 (pencil):
% the homotopy-lifting square A -> X, A -> A x I (cofibration), p : X -> Y,
% h : A x I -> Y, with a diagonal; its dual with a path object A^I, d_0 and
% a cofibration Y -> X; a plain lifting square A -> X, B -> Y. Page 23
% (pencil): two drafts of the square of page 22 — i : X -> X', p : Y -> Y',
% two pairs of parallel maps g_0, g_1 : X -> Y and g'_0, g'_1 : X' -> Y'
% with homotopies k and h between them, and phi : X' -> Y. Page 22 (« e »):
% « We know there exist homotopies k : g_0 = phi i -> g_1 = g, h : g'_0 =
% p phi -> g'_1 = g'. All we have to do is to check that we can choose these
% homotopies compatible with i and p » — trivially so when Y' is the final
% object. A first attempt through a coproduct (X u X) u (X' u X') -> X (x) J
% u X' (x) J is struck; the square X (x) J -> Y, X' (x) J -> Y' is « a priori
% not commutative, but only commutative on X u X -> X (x) J »; the lower half
% of the page, which reduces the problem to a lift k' in a square X u X ->
% Y, j : X u X -> X (x) J (« trivial cofibration »), p (« fibration »), v =
% h o (i (x) J), and ends « pk != v, only pk ~s v », is cancelled by two long
% crossing strokes. Page 24 (« d »): for X_oo the union of a sequence of
% cofibrations X_i -> X_{i+1}, Hom(X_oo, Y) -> lim Hom(X_i, Y) is a
% surjection; to prove injectivity, extend (f, f') from X_oo u X_oo to
% X_oo x Delta(1) step by step, via B'_{i+1} = B_i u_{A_i} A_{i+1} inside
% B_{i+1} = Delta(1) x X_{i+1}; then the question: given Y <= Z_0 -> Z_1 ->
% ..., with an extension Z_i -> Y for every i, is there one on Z_oo? « not
% clear at all », the whole point being that lim pi_0 (?)(Z_i; Y) is non-
% empty — « Ask Ronnie Brown! ».
%
% Notation fixed once for the batch: X (x) J as $X \otimes J$ (J his
% interval object, written J on pp. 22-23; on p. 21 he writes A x I and A^I
% with I); Delta(1) as $\Delta(1)$; X_oo as $X_\infty$; his « ~ » for
% homotopic as $\sim$, and the « ~ » with an « s » above as
% $\overset{s}{\sim}$; the coproduct as $\sqcup$ and the amalgamated sum as
% $\sqcup_{A_i}$. His circled page letters are given in a \note. Greek letters
% are kept in math mode throughout. Inclusions he draws as hooked strokes
% without heads are given as hooked arrows, said once in a note (p. 24).
%
% Readings to check first: « by assumption » and « simplicial » inserted
% above line 1 of p. 22; the letter of the first homotopy there (read k, as
% in the diagrams, but written like an h); the struck word before
% « shortcutting » and « shortcutting » itself; « unhappily »-like word left
% \ill{}; « non » struck before « trivial cofibration »; on p. 24 « surjection »,
% « extend » (« We have to extend step by step given homotopies »), the
% {0}, {1} of the last union, « point », and the underlined overwritten word
% after pi_0 (left \ill{}); on p. 23 the labels of the first draft's arcs.

\documentclass[11pt,a4paper]{article}
% Le préambule commun à toutes les transcriptions du fonds Grothendieck.
%
% Il tient en une idée : **l'incertitude se compose, elle ne se résout pas.**
% Une transcription qui lisse un mot douteux a détruit la seule chose pour
% laquelle on la faisait — un lecteur doit pouvoir savoir, à chaque endroit,
% ce qui était sur le feuillet et ce qui a été ajouté. Les six macros qui
% suivent sont donc l'appareil critique, et non de la décoration.
%
% Les mêmes commandes sont reconnues par `scripts/render.mjs`, qui produit la
% vue de lecture du site. Une macro ajoutée ici doit l'être aussi là-bas,
% faute de quoi le rendu échouera bruyamment — ce qui est voulu : un rendu
% silencieusement faux est pire qu'un rendu absent.

\ProvidesPackage{grothendieck}[2026/08/08 Transcriptions du fonds Grothendieck]

\usepackage{amsmath,amssymb,amsthm}
\usepackage{mathrsfs}
\usepackage{stmaryrd}
% Les diagrammes commutatifs sont partout dans ce fonds : c'est la langue
% même de la géométrie algébrique de Grothendieck, et une transcription qui
% les rendrait en prose ne transcrirait rien.
\usepackage{tikz-cd}
% Français par défaut — c'est la langue des feuillets — mais l'anglais est
% chargé aussi : la traduction et le résumé sont des documents anglais, et la
% typographie française (espaces avant les deux-points) y serait une faute.
% Ces documents basculent par \selectlanguage{english} en tête de corps.
\usepackage[english,french]{babel}
\usepackage{fontspec}
\usepackage{geometry}
\geometry{a4paper,margin=25mm,marginparwidth=28mm}
\usepackage{marginnote}
\usepackage{xcolor}
% ulem, not soul. `soul`'s \st and \ul are built on a character-by-character
% scan that XeLaTeX's font machinery breaks: the document fails with an
% undefined control sequence rather than degrading. ulem does the same two jobs
% and is Unicode-engine safe. `normalem` stops it hijacking \emph, which the
% transcriptions use for Grothendieck's own emphasis.
\usepackage[normalem]{ulem}
\usepackage{hyperref}
\hypersetup{colorlinks=true,linkcolor=black,urlcolor=black,pdfborder={0 0 0}}

% --- Métadonnées du lot -----------------------------------------------------
% Reprises telles quelles par le script de rendu, qui les affiche en tête de
% la vue de lecture. Elles fixent ce que le fichier prétend couvrir : sans
% elles, une transcription flotte sans point d'attache dans un fonds de seize
% mille feuillets.

\newcommand{\folder}[1]{\gdef\grothFolder{#1}}
\newcommand{\batch}[1]{\gdef\grothBatch{#1}}
\newcommand{\leaves}[2]{\gdef\grothFirst{#1}\gdef\grothLast{#2}}
\newcommand{\foldertitle}[1]{\gdef\grothFoldertitle{#1}}
\gdef\grothFolder{?}\gdef\grothBatch{?}\gdef\grothFirst{?}\gdef\grothLast{?}\gdef\grothFoldertitle{}

% --- Le résumé --------------------------------------------------------------
%
% Il ouvre la lecture modernisée, sous le titre « Résumé », et il est composé
% en petit. Ce n'est pas un ornement : c'est le seul passage du document qui
% ne soit pas des mathématiques, et un lecteur doit pouvoir le reconnaître
% avant de l'avoir lu — puis le sauter s'il connaît déjà le sujet, ou s'y
% arrêter s'il ne le connaît pas. Le filet qui le suit marque la reprise du
% corps.

\newenvironment{resume}
  {\small\setlength{\parskip}{0.45em}}
  {\par\medskip\hrule\medskip}

% --- Mots-clés ---------------------------------------------------------------
%
% En anglais, en clôture du résumé de la lecture modernisée : le vocabulaire
% moderne sous lequel chercher ce que les feuillets construisent. Le manifeste
% du site (`npm run manifest`) les extrait pour étiqueter la cote entière —
% cette ligne est la source unique des tags, il n'y a pas de fichier à côté.
\newcommand{\keywords}[1]{\par\smallskip\noindent
  {\footnotesize\textsc{Keywords} --- \textit{#1}}\par}

% --- L'appareil critique ----------------------------------------------------

% Le numéro de feuillet, en marge. C'est l'ancre qui permet de régler un
% désaccord sur une formule en regardant *un* feuillet plutôt qu'une chemise
% de six cents. Le site s'en sert aussi pour tourner les pages du fac-similé
% quand on fait défiler la transcription.
\newcommand{\leaf}[1]{%
  \marginnote{\footnotesize\color{gray!70}\textsf{#1}}%
  \label{leaf:#1}\ignorespaces}

% Illisible : on ne devine pas. Un mot inventé qui se lit comme les autres est
% le pire résultat possible.
\newcommand{\ill}{\textcolor{red!60!black}{[\,\ldots\,]}}

% Lecture douteuse : le mot est donné, et signalé comme incertain.
\newcommand{\uncertain}[1]{\uline{#1}}

% Ajout de l'éditeur — une lettre manquante, un mot sous-entendu.
\newcommand{\add}[1]{\textcolor{blue!50!black}{[#1]}}

% Biffé par Grothendieck lui-même. Ce qu'il a rayé fait partie du document :
% une rature dit souvent où la pensée a changé de direction.
\newcommand{\struck}[1]{\sout{#1}}

% Note du transcripteur, sur l'état matériel du feuillet.
\newcommand{\note}[1]{\par\smallskip{\small\color{gray!60!black}\textit{[#1]}}\par\smallskip}

% Note marginale de Grothendieck, distincte du corps du texte.
\newcommand{\marginal}[1]{\par\smallskip\noindent\hspace{1em}%
  {\small\color{gray!60!black}#1}\par\smallskip}

% --- Titre ------------------------------------------------------------------

\AtBeginDocument{%
  \begin{center}
    {\large\bfseries Fonds Alexandre Grothendieck --- cote n\textsuperscript{o}~\grothFolder}\\[2pt]
    {\itshape\grothFoldertitle}\\[4pt]
    {\small feuillets \grothFirst--\grothLast\ (lot \grothBatch)}\\[2pt]
    {\footnotesize\color{gray!60!black}Transcription établie d'après le fac-similé numérisé
     par l'Université de Montpellier.\\
     Les lectures incertaines sont soulignées, les illisibles notées [\,\ldots\,],
     les ajouts entre crochets.}
  \end{center}
  \vspace{1em}}

\endinput


\folder{107}
\batch{2}
\pages{21}{25}
\dating{[à partir de 1982]}
\watermark{Édition de démonstration}
\foldertitle{Closed models [notes postérieures à PS ?] : notes manuscrites (s.d.)}

\begin{document}

\page{21}

\note{au crayon, dans la marge gauche d'une feuille de listing ; trois
diagrammes, sans texte}

\begin{tikzcd}
A \arrow[r, "\alpha"] \arrow[d, "\text{cof.}"'] & X \arrow[d, "p"] \\
A \times I \arrow[r, "h"'] \arrow[ur] & Y
\end{tikzcd}

\note{à droite de ce carré, deux petits croquis : une flèche verticale
prise dans une lentille (deux arcs), deux fois, le second arc supérieur
portant une pointe — vraisemblablement le dessin d'une homotopie entre deux
flèches}

\begin{tikzcd}
A & X \arrow[l] \\
A^{I} \arrow[u, "\partial_0"] & Y \arrow[u, "\text{cofib.}"'] \arrow[l, Rightarrow]
\end{tikzcd}

\note{à gauche de $A$, « fib. » ; à gauche de la flèche $\partial_0$, un mot
biffé, \struck{\ill{}}. La flèche du bas, $Y \to A^I$, est tracée en trait
double et appuyé}

\note{plus bas, deux signes isolés : un trait oblique et, de lecture
douteuse, $Y_{\times}$}

\begin{tikzcd}
A \arrow[r] \arrow[d] & X \arrow[d] \\
B \arrow[r] \arrow[ur] & Y
\end{tikzcd}

\page{22}

\note{en haut, cerclé, « e » : sa lettre de page. Dans l'ordre des
archivistes, cette page précède la page 24, marquée « d »}

\begin{tikzcd}[column sep=6em, row sep=6em]
X \arrow[r, bend left=35, "g_0 = \varphi i"] \arrow[r, bend right=35, "g_1 = g"'] \arrow[d, "i"'] & Y \arrow[d, "p"] \\
X' \arrow[r, bend left=35, "g'_0 = p\varphi"] \arrow[r, bend right=35, "g'_1 = g'"'] \arrow[ur, bend left=15, "\theta_0 = \varphi"] \arrow[ur, dashed, bend right=15, "\theta_1 = \overline{\varphi}"'] & Y'
\end{tikzcd}

\note{entre les deux arcs du haut, une flèche double oblique marquée $k$ ;
entre ceux du bas, une flèche marquée $h$ (lecture \uncertain{$h$}). La
diagonale $\theta_0 = \varphi$ est pleine, $\theta_1 = \overline{\varphi}$
en tirets ; l'indice de $\theta_0$ est surchargé}

We know \uncertain{by assumption} there exist \uncertain{simplicial}
homotopies\note{« by assumption » et « simplicial » sont ajoutés au-dessus
de la ligne, avec des traits de renvoi}
\[
k : g_0 = \varphi i \longrightarrow g_1 = g
\]
\ill{}
\[
h : g'_0 = p\varphi \longrightarrow g'_1 = g' .
\]
\note{lecture douteuse : la lettre de la première homotopie est tracée comme un $h$ ; les
diagrammes de la page et de la page 23 la nomment $k$}

All we have to do, is to check that we can choose these homotopies
\emph{compatible} with $i$ and $p$\note{trois traits verticaux dans la
marge, devant la parenthèse qui suit}
(and in case $Y'$ is the final object, we are through!)\note{« we are
through! » est écrit au-dessus de \struck{Consider the}, biffé}

\note{à gauche de ce texte, un premier diagramme, barré par un grand trait
courbe ; plusieurs de ses termes ($X \sqcup X$, le premier $X$ de la
deuxième ligne, le premier $\sqcup$, le $Y$ du bas) sont en outre barrés
chacun d'un trait oblique}

\begin{tikzcd}[column sep=small, nodes={font=\small}]
(X \sqcup X) \arrow[d, "\text{cofib?}"'] & \sqcup & (X' \sqcup X') \arrow[d, "\text{cofib?}"'] & \\
X \otimes J & \sqcup \arrow[d] & X' \otimes J \arrow[r] & (Y \otimes J) \\
Y & \sqcup & Y' &
\end{tikzcd}

\note{le terme d'arrivée de la flèche horizontale est de lecture douteuse :
\uncertain{$(Y \otimes J)$} ; il est suivi de \struck{$\sqcup\ \overline{\theta}$}, biffé}

diagram\note{« diagram » n'est pas biffé ; il reste seul sous le « Consider
the » biffé, au-dessus du carré qui suit}

\begin{tikzcd}
X \otimes J \arrow[r, "k"] \arrow[d, "i \otimes J"'] & Y \arrow[d, "p"] \\
X' \otimes J \arrow[r, "h"'] & Y'
\end{tikzcd}

\note{une lentille (deux arcs) va de $X \otimes J$ à $Y'$ à travers le
carré}

which a priori is \emph{not} commutative, but only commutative
\struck{after} on $X \sqcup X \longrightarrow X \otimes J$.

\note{à partir d'ici, toute la moitié inférieure de la page est annulée par
deux longs traits croisés ; elle est transcrite telle qu'elle se lit. À
gauche de la phrase précédente, une première esquisse du même carré,
barrée :}

\begin{tikzcd}
X \otimes J \arrow[r, "k"] \arrow[d, "i \otimes J"'] & Y \arrow[d, "p"] \\
X' \otimes J \arrow[r, "h"'] & Y'
\end{tikzcd}

\begin{tikzcd}[column sep=large]
X \sqcup X \arrow[d] \arrow[rd, bend left=20, "{u = (g_0,\ g_1)}"] & \\
X \otimes J \arrow[r, dashed, "k'?"] \arrow[d, "i \otimes J"'] & Y \arrow[d, "p"] \\
X' \otimes J \arrow[r, "h"'] & Y'
\end{tikzcd}

this diagram is commutative and we look for $k'$ which keeps it
commutative.

We rewrite the problem, \struck{\ill{}} \uncertain{shortcutting}
$X' \otimes J$ by the arrow
\[
h \circ (i \otimes J) = v .
\]
\note{« shortcutting » est écrit au-dessus du mot biffé}

\begin{tikzcd}[column sep=large]
X \sqcup X \arrow[r, "{u = (g_0,\ g_1)}"] \arrow[d, "j"'] & Y \arrow[d, "p\ \text{fibration}"] \\
X \otimes J \arrow[r, "v"'] \arrow[ur, dashed, "k'?"] & Y'
\end{tikzcd}

\note{à gauche de la flèche $j$ : \uncertain{\struck{non}} trivial
cofibration ; plus à gauche encore, un signe biffé, \struck{\ill{}}}

We know there is a \struck{and we are \ill{}} ! $k$ s.th.\ $kj = u$, but
\ill{} $pk \neq v$, only $pk \overset{s}{\sim} v$\note{le mot avant
$pk \neq v$ est surchargé et souligné}

(But in case $Y'$ is final object we are evidently through!)

\page{23}

\note{au crayon, sur une feuille de listing tenue de travers : deux
esquisses du premier diagramme de la page 22, sans texte. Le listing est
daté 82155, soit début juin 1982}

\begin{tikzcd}[column sep=6em, row sep=4em]
X \arrow[r, bend left=35, "g_0 = \varphi i"] \arrow[r, bend right=35, "g_1 = \varphi i"'] \arrow[d] & Y \arrow[d] \\
X' \arrow[r, bend left=35, "g'_0 = g'"] \arrow[r, bend right=35, "g'_1 = p\varphi"'] & Y'
\end{tikzcd}

\note{première esquisse : les étiquettes des arcs ne sont pas celles de la
page 22 ; on les donne comme elles se lisent, et celles de $g_1$ et de
$g'_0$ sont douteuses}

\begin{tikzcd}[column sep=6em, row sep=6em]
X \arrow[r, bend left=35, "g_0 = \varphi i"] \arrow[r, bend right=35, "g_1 = g"'] \arrow[d, "i"'] & Y \arrow[d, "p"] \\
X' \arrow[r, bend left=35, "g'_0 = p\varphi"] \arrow[r, bend right=35, "g'_1 = g'"'] \arrow[ur, "\varphi" description] & Y'
\end{tikzcd}

\note{seconde esquisse : entre les arcs du haut une flèche marquée $k$,
entre ceux du bas une flèche marquée $h$ ; le $Y$ du haut est surchargé}

\page{24}

\note{en haut, cerclé, « d » : sa lettre de page}

Returning to our initial pb, we get \struck{\ill{}}\note{un mot court,
biffé, en tête de la formule}
\[
\operatorname{Hom}(X_\infty, Y) \longrightarrow
\varprojlim \operatorname{Hom}(X_i, Y)
\]
\emph{\uncertain{surjection}} (using strongly that the transition
morphisms $X_i \to X_{i+1}$ are cofibrations). To prove it injective, let
\[
f, f' : X_\infty \longrightarrow Y, \qquad
f_i = f \,|\, X_i, \quad f'_i = f' \,|\, X_i ;
\]
\uncertain{we've} to prove that $f_i \sim f'_i$ $\forall i$, implies
$f \sim f'$.

Consider
\[
\underbrace{X_\infty \times \Delta(1)}_{B}
\ \underset{\text{cofib.}}{\supset}\
\underbrace{X_\infty \sqcup X_\infty}_{A} ,
\]
with $(f, f') : A \to Y$; we want to extend to $X_\infty \times \Delta(1)$,
and we know it is possible upon $X_i \times \Delta(1)$ $\forall i$.
\note{la flèche $(f,f')$ est tracée de $A$ vers $Y$, sous la formule}

\[
\begin{array}{ccc}
B & \supset & A \\
\vdots & & \vdots \\
B_{i+1} & \supset & A_{i+1} \\
\text{cofib}\ \cup & & \cup\ \text{cofib} \\
B_i & \underset{\text{cofib.}}{\supset} & A_i
\end{array}
\]

We have to \uncertain{extend} step by step given homotopies. The problem is
to construct $h_{i+1}$ extending $h'_{i+1}$.

\begin{tikzcd}[column sep=4em, row sep=3.5em]
A_i \arrow[r, hook] \arrow[d, hook] \arrow[ddr, bend right=70, "\varphi_i"'] & B_i \arrow[d, hook] & \\
A_{i+1} \arrow[r, hook] \arrow[dr, "\varphi_{i+1}"'] & B'_{i+1} \arrow[r, hook] \arrow[d, "h'_{i+1}"] & B_{i+1} \arrow[dl, dashed, "h_{i+1}?"] \\
 & Y &
\end{tikzcd}
\[
B'_{i+1} \overset{\text{def}}{=} B_i \sqcup_{A_i} A_{i+1}
\]
\note{les lettres $\varphi_i$, $\varphi_{i+1}$ sont de lecture douteuse.
Les inclusions sont tracées en traits à crochet, sans pointe ; un
trait vertical relie $B_i$ au milieu du segment $A_{i+1}$ — $B'_{i+1}$. Sur
la flèche en tirets $B_{i+1} \to Y$ est d'abord tracée une flèche pleine,
raturée. À gauche du diagramme, un trait isolé}

\[
\begin{array}{c}
\Delta(1) \times X_{i+1} \\
\cup \\
\bigl(\Delta(1) \times X_i\bigr) \cup \bigl(\{0\} \times X_{i+1}\bigr) \cup \bigl(\{1\} \times X_{i+1}\bigr)
\end{array}
\]
\note{les $\{0\}$ et $\{1\}$ sont surchargés et de lecture douteuse}

\[
Y \Leftarrow Z_0 \hookrightarrow Z_1 \hookrightarrow Z_2 \cdots
\]
Assume $\forall i$, $\exists$ extension $Z_i \to Y$. Is there an extension
$Z_\infty \to Y$?

\emph{not clear at all}, the \struck{\ill{}} the whole \uncertain{point} is
to see that
$\varprojlim_i \pi_0$ \ill{} $(Z_i ; Y) \neq \emptyset$ ---\note{« the whole » est
ajouté au-dessus du mot biffé, avec un trait de renvoi ; le mot après
$\pi_0$ est souligné et surchargé}

Ask Ronnie Brown!

\end{document}
